\section{Eval 性能与真实能力的鸿沟}

\subsection{模型的精神分裂}

Ilya 指出当前模型存在一个令人困惑的现象：它们在 eval 上表现极好，但在实际使用中仍然犯低级错误。他用 vibe coding 的例子说明：模型修了一个 bug，引入了第二个 bug；你告诉它第二个 bug，它在修复时又恢复了第一个 bug——可以在两个 bug 之间无限循环。

\begin{importantbox}{两种可能的解释}
\begin{enumerate}
    \item \textbf{RL 使模型过度单一聚焦}：RL 训练可能使模型在某些方面变得"too single-minded"，虽然在其他方面也增强了能力，但基本判断力可能受损
    \item \textbf{RL 环境设计受 eval 启发}：研究团队在设计 RL 训练环境时不自觉地从 eval 中汲取灵感（"我希望模型在这个 eval 上表现好，什么样的 RL 训练能帮到这里？"），导致模型实际上在"teaching to the test"
\end{enumerate}
\end{importantbox}

\subsection{10000 小时 vs 100 小时的类比}

\begin{importantbox}{竞赛选手 vs 天才选手}
假设两个学生学习竞赛编程：
\begin{itemize}
    \item 学生 A 练习了 10,000 小时，解题无数，记住所有证明技巧——成为顶尖选手
    \item 学生 B 只练习了 100 小时，但同样表现出色
\end{itemize}
谁的职业发展前景更好？\textbf{第二个}。因为学生 B 的能力来自更深层的泛化能力，而非特定领域的过度训练。当前的 AI 模型更像学生 A——甚至更极端，因为我们不仅使用了所有竞赛题，还做了数据增广来生成更多题目。
\end{importantbox}

\begin{warningbox}{真正的 Reward Hacking}
Ilya 暗示：真正的 reward hacking 不是模型在做——而是\textbf{人类研究者}在做。他们过度关注 eval，不自觉地将 RL 训练环境对齐到 eval 指标，造成了 eval 性能与真实能力之间的系统性偏差。
\end{warningbox}

\subsection{本章小结}

Eval 性能与真实世界能力之间存在显著鸿沟。原因可能是 RL 训练过度聚焦于可验证域，以及人类研究者不自觉地"teaching to the test"。模型更像过度训练的竞赛选手，而非真正的通才。


